\section{Pruebas y Verificación}
\label{sec:pruebas}

\subsection{Los siete casos mínimos de funcionamiento}

La Tabla~\ref{tab:casos} resume los siete casos mínimos que exige la sección
"Pruebas Mínimas de Funcionamiento" del enunciado, cada uno invocable de
forma aislada desde \texttt{scripts/casos\_enunciado/} y en conjunto vía
\texttt{make casos-enunciado}. Los comandos exactos, los \textit{fixtures}
usados y la implementación de cada caso están documentados en el README del
repositorio de código y en \texttt{milestone10\_notas\_tecnicas.md} del
repositorio de conocimiento.

\begin{table}[htbp]
\caption{Los 7 casos mínimos: comando, esperado, obtenido, conclusión}
\label{tab:casos}
\renewcommand{\arraystretch}{1.25}
\begin{tabular}{p{1.55cm}p{2.7cm}p{3.3cm}}
\toprule
\textbf{Caso} & \textbf{Esperado} & \textbf{Obtenido} \\
\midrule
1. Validación &
  Rechazo con código $\neq 0$ y sin ejecución parcial &
  8/8 verificaciones cumplen (los 4 escenarios del enunciado + evidencia de sin ejecución parcial) \\
\midrule
2. Reproducibilidad &
  Dos corridas idénticas; las 5 tareas terminan &
  Logs idénticos byte a byte; 5/5 terminan \\
\midrule
3. Igualdad &
  Share medio $\approx 0.20$, con desviación y error &
  Error absoluto medio $0.000477$ (tolerancia $0.02$), 30 semillas \\
\midrule
4. Proporcionalidad &
  Shares objetivo $1/15 \ldots 5/15$; error $\leq 0.02$ &
  Error absoluto medio $0.001104$; shares 1/15\ldots5/15 confirmados \\
\midrule
5. Terminación &
  Ninguna FINISHED vuelve a ganar; suma correcta; joins completos &
  15/15 verificaciones cumplen \\
\midrule
6. Modos &
  Mismo trabajo final y $\pi$ entre modos &
  7/7 verificaciones cumplen; $\pi$ idéntico bit a bit \\
\midrule
7. Estrés &
  Sin deadlock, fuga, acceso inválido, \textit{data race} ni UB &
  12/12 en las tres variantes de \textit{sanitizer} (ver \S\ref{sec:sanitizers}) \\
\bottomrule
\end{tabular}
\end{table}

Los Casos 1, 5 y 6 reportan únicamente el escenario que nombra el enunciado
--- no toda la suite de ingeniería heredada de los milestones donde se
originaron (Milestone 1/8 para la validación de CLI/CSV, Milestone 5 para el
bucle del planificador, Milestone 6 para los modos de ejecución). Esa
cobertura adicional (52 verificaciones en total: 20 de validación extra de
CLI/CSV, 20 del bucle del planificador, 39 de los modos de ejecución más 9
de la validación de \texttt{--yield-config}) se mantiene como parte del
ciclo normal de desarrollo (\texttt{make test}), separada de la evidencia
mínima que exige el enunciado.

\subsection{Evidencia de \textit{sanitizers}}
\label{sec:sanitizers}

El Caso 7 (Estrés: 25 tareas, parámetros variados, 5 semillas en ambos modos
más los casos borde de $Q=1$ y \texttt{--max-dispatches}) se verificó en tres
variantes, cada una en el entorno donde certifica algo distinto:

\begin{itemize}
\item \textbf{AddressSanitizer + UndefinedBehaviorSanitizer}, dentro de
  Docker (Ubuntu 24.04 x86\_64, el entorno de referencia del enunciado):
  12/12, sin hallazgos.
\item \textbf{AddressSanitizer + LeakSanitizer}, dentro de Docker: 12/12,
  sin fugas de memoria.
\item \textbf{ThreadSanitizer}, fuera de Docker, directo en el host: 12/12,
  sin \textit{data races}.
\end{itemize}

Estas tres variantes no corren en el mismo entorno por dos limitaciones
conocidas, no atribuibles al código del proyecto: \textbf{LeakSanitizer no
reporta fugas de forma confiable fuera de un Linux real} (un
\texttt{malloc} deliberado sin \texttt{free}, compilado con
\texttt{-fsanitize=address} y corrido fuera de Docker, termina con
\texttt{exit=0} sin reportar nada; el mismo binario dentro de Docker sí lo
detecta); y \textbf{ThreadSanitizer falla dentro de Docker en hosts
ARM64} (\texttt{FATAL: ThreadSanitizer: unexpected memory mapping}) por un
problema conocido de su instrumentación bajo la emulación x86\_64 de
Rosetta/QEMU, no relacionado con el código. Ninguna de las dos es un falso
positivo del análisis en sí --- son limitaciones del entorno de ejecución
del \textit{sanitizer}, mitigadas corriendo cada uno donde efectivamente
funciona.
